\part{Introduction}
\section{Introduction}

%%%%%%%%%%%%%%%%%%%%%%%%%%%%%%%%%%%%%%%%%%%%%%%%%%%%%%%%%%%%%%%%%%%%%%%%%%%%%%%%%%%%%%%%%%%%%%%%%%%%%%%%%%%%%%%%%%%%%%%%%%%
%%%%%%%%%%%%%%%%%%%%%%%%%%%%%%%%%%%%%%%%%%%%%%%%%%%%%%%%%%%%%%%%%%%%%%%%%%%%%%%%%%%%%%%%%%%%%%%%%%%%%%%%%%%%%%%%%%%%%%%%%%%

\begin{frame}
	\partpage
	\centering
\end{frame}

%%%%%%%%%%%%%%%%%%%%%%%%%%%%%%%%%%%%%%%%%%%%%%%%%%%%%%%%%%%%%%%%%%%%%%%%%%%%%%%%%%%%%%%%%%%%%%%%%%%%%%%%%%%%%%%%%%%%%%%%%%%
%%%%%%%%%%%%%%%%%%%%%%%%%%%%%%%%%%%%%%%%%%%%%%%%%%%%%%%%%%%%%%%%%%%%%%%%%%%%%%%%%%%%%%%%%%%%%%%%%%%%%%%%%%%%%%%%%%%%%%%%%%%

\begin{frame}
	\frametitle{A pan-genomic algorithm}
	\centering
	
	From the greek word \textit{pan} (everything)
	
	\medskip
	
	Given a series of complete genomes
	
	\medskip

	\includegraphics[height=0.4cm]{images/gen1}

  \medskip
  	
	\includegraphics[height=0.4cm]{images/gen2}
	
	\pause
	
	\medskip

  What we want to see:
  
  \begin{itemize}
    \item {\color{orange}Similarities}
    \item {\color{purple}Differences}
    \item {\color{green}Duplications}
  \end{itemize}	
	
	\medskip
	
	\includegraphics[height=1cm]{images/gen3}
	
\end{frame}

%%%%%%%%%%%%%%%%%%%%%%%%%%%%%%%%%%%%%%%%%%%%%%%%%%%%%%%%%%%%%%%%%%%%%%%%%%%%%%%%%%%%%%%%%%%%%%%%%%%%%%%%%%%%%%%%%%%%%%%%%%%
%%%%%%%%%%%%%%%%%%%%%%%%%%%%%%%%%%%%%%%%%%%%%%%%%%%%%%%%%%%%%%%%%%%%%%%%%%%%%%%%%%%%%%%%%%%%%%%%%%%%%%%%%%%%%%%%%%%%%%%%%%%

\begin{frame}
	\frametitle{de Bruijn graph}
	\centering
	  
  Sequence \#1: {\color{red}TGACGTC} ($2$-mers: {\color{red} TG GA AC CG GT TC})
                            
  \medskip

	\includegraphics[height=0.7cm]{images/debruin-0}

  \medskip
  
  Sequence \#2: {\color{blue}TGACTTC} ($2$-mers: {\color{blue} TG GA AC CT TT TC})
	
	\medskip
	
	\includegraphics[height=0.7cm]{images/debruin-1}
	
	\pause

	\medskip
	
	Find common vertices

	\medskip
  \includegraphics[height=1.5cm]{images/debruin-2}

	\pause
	                           
  \medskip
  
  Merge them
	\medskip

	\includegraphics[height=1.3cm]{images/debruin-3}


\end{frame}

%%%%%%%%%%%%%%%%%%%%%%%%%%%%%%%%%%%%%%%%%%%%%%%%%%%%%%%%%%%%%%%%%%%%%%%%%%%%%%%%%%%%%%%%%%%%%%%%%%%%%%%%%%%%%%%%%%%%%%%%%%%
%%%%%%%%%%%%%%%%%%%%%%%%%%%%%%%%%%%%%%%%%%%%%%%%%%%%%%%%%%%%%%%%%%%%%%%%%%%%%%%%%%%%%%%%%%%%%%%%%%%%%%%%%%%%%%%%%%%%%%%%%%%

\begin{frame}
	\frametitle{Compacted de Bruijn graph}
	\centering
	
  To compact the de Bruijn graph we need to find the junctions
  
  \medskip
  
  A vertex is a junction if :
  
  Is the initial or final part of some sequence\\
  Has more than 1 in-edge or more than 1 out-edge\\
	
	\medskip

	\includegraphics[height=2cm]{images/debruin-3a}

  \pause

	\medskip

  Remove simple paths by connect junctions  	

	\medskip

	\includegraphics[height=1.5cm]{images/debruin-4}
	
\end{frame}

%%%%%%%%%%%%%%%%%%%%%%%%%%%%%%%%%%%%%%%%%%%%%%%%%%%%%%%%%%%%%%%%%%%%%%%%%%%%%%%%%%%%%%%%%%%%%%%%%%%%%%%%%%%%%%%%%%%%%%%%%%%
%%%%%%%%%%%%%%%%%%%%%%%%%%%%%%%%%%%%%%%%%%%%%%%%%%%%%%%%%%%%%%%%%%%%%%%%%%%%%%%%%%%%%%%%%%%%%%%%%%%%%%%%%%%%%%%%%%%%%%%%%%%

%\begin{frame}
%	\frametitle{Junctions}
%	\centering
%\end{frame}

%%%%%%%%%%%%%%%%%%%%%%%%%%%%%%%%%%%%%%%%%%%%%%%%%%%%%%%%%%%%%%%%%%%%%%%%%%%%%%%%%%%%%%%%%%%%%%%%%%%%%%%%%%%%%%%%%%%%%%%%%%%
%%%%%%%%%%%%%%%%%%%%%%%%%%%%%%%%%%%%%%%%%%%%%%%%%%%%%%%%%%%%%%%%%%%%%%%%%%%%%%%%%%%%%%%%%%%%%%%%%%%%%%%%%%%%%%%%%%%%%%%%%%%

\begin{frame}
	\frametitle{The problem}
	\centering
	
	How to construct the compacted de Bruijn graph,\\ from many complete genomes,\\ avoiding the ordinary graph construction?
	
	\pause
	
	\medskip
	
	$\downarrow$
	
	\medskip

  
  \textbf{Which $k$-mers are junctions?}
  
  \medskip
  
  +
  
  \medskip
  
  How to construct the compacted de Bruijn graph given the junctions?
  
	
\end{frame}
